\chapter{Regression Methods}
\label{chap:regression-methods}

\section{Regression Framework}
\label{sec:reg-framework}

\subsection{Linear Regression}
The linear regression model is the foundation for analyzing the relationship between a continuous outcome $Y$ and a set of predictors $\mathbf{X} = \{X_1, X_2, \dots, X_p\}$ \cite{mccullagh1989}. The model is defined as:
\begin{equation}
    Y = \beta_0 + \sum_{j=1}^{p} \beta_j X_j + \epsilon, \quad \epsilon \sim \mathcal{N}(0, \sigma^2)
\end{equation}
where $\beta_j$ are the regression coefficients and $\epsilon$ represents the error term.

\subsection{Logistic Regression}
For binary outcomes $Y \in \{0, 1\}$, we employ the logistic regression model \cite{hosmer2013}:
\begin{equation}
    \text{logit}(P(Y=1|\mathbf{X})) = \ln\left(\frac{P(Y=1|\mathbf{X})}{1-P(Y=1|\mathbf{X})}\right) = \beta_0 + \sum_{j=1}^{p} \beta_j X_j
\end{equation}

\subsection{Inferential Quantities of Interest}
Our evaluation focuses on the preservation of the estimated coefficients $\hat{\beta}_j$ and their associated 95\% confidence intervals (CIs). If synthetic data is to be used as a proxy for real data, the synthetic coefficients $\hat{\beta}_j^{(S)}$ should be unbiased estimates of the original coefficients $\hat{\beta}_j^{(R)}$, and their CIs should overlap substantially.

\section{Data Framework}
\label{sec:reg-data-framework}

\subsection{Original Datasets and Synthetic Datasets}
We generate Original Datasets (OD) through a controlled simulation process. From each OD, we create one corresponding Synthetic Dataset (SD) using a specified generator.
More explicitly, the predictor variables are first generated from a multivariate Normal distribution, and the response variable is then generated according to the relevant downstream model: a Gaussian linear model for linear regression or a Bernoulli model with logistic link for logistic regression. This OD-to-SD workflow is the one implemented in the experiment configuration reported in Appendix~\ref{sec:app-json}.

\subsection{General Comparison Strategy: OD versus SD}
The utility of the SD is measured by comparing the results of the same regression model fitted to both the OD and the SD. This ``one-to-one'' comparison ensures that any observed degradation is attributable solely to the synthesis process.

\subsection{Role of Repeated Simulation Scenarios}
To ensure statistical robustness, each unique scenario (combination of factors) is repeated $M = 1,000$ times. This allows us to calculate the average performance and the variability of the evaluation metrics across independent runs.

\section{Data Generation Mechanisms}
\label{sec:reg-generation-mechanisms}

\subsection{Simulation Factors}
The simulation space is defined by the following factors:
\begin{equation}
N \in \{100, 500, 1000\}, \qquad
p \in \{2, 5, 10\}, \qquad
\rho \in \{0.0, 0.3, 0.6\}
\end{equation}
\begin{equation}
\sigma^2 \in \{0.5, 1.0, 2.0\}
\end{equation}
where $\sigma^2$ applies to the linear-regression scenarios, while the logistic-regression scenarios vary the outcome prevalence through the intercept term.

\subsection{Sample Size, Number of Predictors, and Variable Type}
The design includes the following sample-size, dimensionality, and predictor-type structure:
\begin{equation}
N \in \{100, 500, 1000\}, \qquad p \in \{2, 5, 10\}
\end{equation}
\begin{equation}
\text{VarType} \in \{\text{continuous}, \text{binary}\}
\end{equation}

\subsection{Correlation Structure and Error Variability}
The continuous predictor block is generated from a multivariate Normal distribution:
\begin{equation}
\mathbf{X} \sim \mathcal{N}_p(\mathbf{0}, \Sigma)
\end{equation}
with covariance structure
\begin{equation}
\Sigma_{ii} = 1, \qquad \Sigma_{ij} = \rho \quad (i \neq j), \qquad \rho \in \{0.0, 0.3, 0.6\}
\end{equation}
For the binary scenarios, the same latent dependence structure is used and then discretized into binary predictors. For linear models, the error variance is set as
\begin{equation}
\sigma^2 \in \{0.5, 1.0, 2.0\}
\end{equation}

\subsection{Regression Coefficients and Outcome Generation}
The coefficients are fixed as
\begin{equation}
\beta_j = (-1)^j, \qquad j = 1, \dots, p
\end{equation}
For the linear-regression scenarios, the response is generated as
\begin{equation}
Y = \beta_0 + \mathbf{X}\beta + \epsilon, \qquad \epsilon \sim \mathcal{N}(0, \sigma^2)
\end{equation}
For the logistic-regression scenarios, the response is generated as
\begin{equation}
Y \sim \operatorname{Bernoulli}(\pi), \qquad \log\left(\frac{\pi}{1-\pi}\right) = \beta_0 + \mathbf{X}\beta
\end{equation}
The same data-generation logic is encoded in the regression configuration summarized in Appendix~\ref{sec:app-json}.

\section{Synthetic Data Generators}
\label{sec:reg-generators}

\subsection{Generators for Predictors}
We evaluate two primary approaches for synthesizing the predictors $\mathbf{X}$:
\begin{itemize}
    \item \textbf{CART:} Non-parametric classification and regression trees.
    \item \textbf{Parametric:}
    \begin{itemize}
        \item \textbf{Normal:} for continuous predictors.
        \item \textbf{Logistic:} for binary predictors.
    \end{itemize}
\end{itemize}

\subsection{Generators for Outcomes}
The outcome $Y$ is synthesized last, conditioned on all previously generated predictors $\mathbf{X}$, using the same family of methods (CART or Parametric).

\subsection{Combined Generation Arms}
The study evaluates ``homogeneous'' arms where the same method family is used for both predictors and outcomes (e.g., CART-X + CART-Y vs. Norm-X + Norm-Y).

\subsection{Conceptual Differences between the Generators}
Parametric generators assume the underlying data follows a specific distribution, whereas CART generators are model-agnostic and can capture non-linearities and interactions, albeit at the risk of overfitting or introducing ``blocky'' artifacts.

\section{Experimental Design}
\label{sec:reg-experimental-design}

\subsection{Scenario Grid}
The full experimental design comprises 1,800+ scenarios, resulting from the Cartesian product of all simulation factors.

\subsection{Number of Repetitions and Successful Fits}
Each scenario is executed until $M=1,000$ successful fits are obtained. Scenarios that fail to converge (common in logistic regression with small $N$ and high $p$) are flagged but excluded from the final metric calculation to ensure comparability.

\subsection{Computational Workflow}
The workflow is implemented in R (using \texttt{synthpop}) for data generation and synthesis, and Python for the evaluation and visualization of results.

\subsection{Quality-Control and Fit-Validation Rules}
We apply strict quality control to ensure that the original data is ``well-behaved'' (no perfect collinearity) before synthesis, ensuring that any failures in the SD are truly reflective of synthetic degradation.

\section{Evaluation Strategy}
\label{sec:reg-evaluation}

A summary of the metrics used to assess inferential utility is presented in Table~\ref{tab:regression_metrics}.

\begin{table}[H]
  \centering
  \caption[Summary of Regression Evaluation Metrics]{Summary of the inferential metrics used to evaluate regression utility.}
  \label{tab:regression_metrics}
  \renewcommand{\arraystretch}{1.3}
  \begin{tabular}{@{}lp{8cm}l@{}}
    \toprule
    \textbf{Metric} & \textbf{Description} & \textbf{Ideal Value} \\ \midrule
    \textbf{CIO} & Confidence Interval Overlap: overlap between original and synthetic 95\% CIs. & $1.0$ (100\%) \\
    \textbf{Bias Ratio} & Ratio of synthetic coefficient bias to the original standard error. & $0.0$ \\
    \textbf{CI Width Ratio} & Ratio of synthetic CI width to original CI width. & $1.0$ \\ \bottomrule
  \end{tabular}
  \source{Compiled by the author}
\end{table}

\subsection{Main Metric: CI Overlap}
The primary measure of inferential utility is the Confidence Interval Overlap (CIO). It is first computed coefficient by coefficient and then averaged over all predictors in the fitted model. For coefficient $j$, we define:
\begin{equation}
    \text{CIO}_j = \frac{1}{2} \left( \frac{\min(U_{R,j}, U_{S,j}) - \max(L_{R,j}, L_{S,j})}{U_{R,j} - L_{R,j}} + \frac{\min(U_{R,j}, U_{S,j}) - \max(L_{R,j}, L_{S,j})}{U_{S,j} - L_{S,j}} \right)
\end{equation}
where $[L_{R,j}, U_{R,j}]$ and $[L_{S,j}, U_{S,j}]$ are the 95\% confidence intervals for coefficient $j$ in the original and synthetic data, respectively. The scenario-level CIO is then defined as
\begin{equation}
    \text{CIO} = \frac{1}{p} \sum_{j=1}^{p} \text{CIO}_j
\end{equation}
That is, the metric directly compares the inferential intervals obtained from the original fit and the corresponding synthetic fit, coefficient by coefficient, and then averages them over the full model.

\subsection{Supporting Metrics: Bias and CI Width}
To understand the \emph{source} of CIO degradation, we measure the Bias Ratio and the CI Width Ratio, again coefficient by coefficient and then averaged over the model.

For coefficient $j$, the Bias Ratio is defined as
\begin{equation}
    B_{R,j} = \frac{|\hat{\beta}^{(S)}_j - \hat{\beta}^{(R)}_j|}{\operatorname{SE}(\hat{\beta}^{(R)}_j)}
\end{equation}
and the model-level Bias Ratio is
\begin{equation}
    B_R = \frac{1}{p} \sum_{j=1}^{p} B_{R,j}
\end{equation}

For coefficient $j$, the CI Width Ratio is defined as
\begin{equation}
    W_{R,j} = \frac{\operatorname{Width}^{(S)}_j}{\operatorname{Width}^{(R)}_j}
\end{equation}
and the model-level CI Width Ratio is
\begin{equation}
    W_R = \frac{1}{p} \sum_{j=1}^{p} W_{R,j}
\end{equation}

Thus, all three inferential metrics are computed per predictor and then averaged over the $p$ coefficients in the model.

\subsection{Univariate Analysis Strategy}
We first analyze each factor ($N, p, \rho, \dots$) independently by averaging the metrics across all other factors.

\subsection{Multivariate Analysis Strategy}
Finally, we use correlation heatmaps and multivariate plots to study the interaction between the synthesis methods and the data characteristics themselves, namely how performance changes jointly with factors such as $N$, $p$, $\rho$, and the predictor type.
